\subsection{完备可赋范的实或复代数中的全纯函数演算}

设 E 是一个实拓扑向量空间。拓扑向量空间 \(C \otimes E\)，即 E 的复化（EVT, II, p.~65），记作 \(\mathrm{E}_{(\mathbf{C})}\)，并且把 E 等同于 \(\mathrm{E}_{(\mathbf{C})}\) 的实拓扑向量子空间，所借助的映射是 \(x \mapsto 1 \otimes x\)。

命题 18. --- 复拓扑向量空间 \(E_{(\mathbf{C})}\) 是可赋范的（相应地，完备的），当且仅当 E 是可赋范的（相应地，完备的）。

\(\mathrm{E}_{(\mathbf{C})}\) 的底实拓扑向量空间同构于 \(E \times E\)。于是 \(\mathrm{E}_{(\mathbf{C})}\) 完备当且仅当 E 完备，且若 \(\mathrm{E}_{(\mathbf{C})}\) 可赋范，则 E 可赋范。

反之，设 E 可赋范。设 p 是定义 E 的拓扑的范数，B 是 p 的单位球。存在 \(\mathrm{E}_{(\mathbf{C})}\) 中 0 的闭均衡邻域 V，包含在 \(B+iB\) 中（EVT, II, p.~66）。于是集合 \(\lambda V\)（其中 \(\lambda\) 取遍 \(R_{+}^{*}\)）构成 \(\mathrm{E}_{(\mathbf{C})}\) 中 0 的基本邻域系。V 的规范函数是 \(\mathrm{E}_{(\mathbf{C})}\) 上的一个范数，它定义 \(\mathrm{E}_{(\mathbf{C})}\) 的拓扑，因而 \(\mathrm{E}_{(\mathbf{C})}\) 可赋范。

注记. --- 设 E 与 F 是 K 上的可赋范拓扑向量空间。从 E 到 F 中的连续线性映射的向量空间 \(\mathcal{L}(\mathrm{E};\mathrm{F})\)，赋以有界收敛拓扑，是一个可赋范拓扑向量空间（EVT, III, p.~14）。

设 E 与 F 是 R 上的可赋范拓扑向量空间。C-线性映射 \(\varphi\colon\mathcal{L}(\mathrm{E};\mathrm{F})_{(\mathbf{C})}\to\mathcal{L}(\mathrm{E}_{(\mathbf{C})};\mathrm{F}_{(\mathbf{C})})\)（由 \(\varphi(\lambda\otimes u)=\lambda u_{(\mathbf{C})}\) 定义）是复拓扑向量空间之间的同构。特别地，\(\mathrm{E}_{(\mathbf{C})}\) 的对偶等同于 E 的对偶的复化，赋范代数 \(\mathcal{L}(\mathrm{E}_{(\mathbf{C})})\) 等同于赋范代数 \(\mathcal{L}(\mathrm{E})\) 的复化。

设 S 是 C 中在复共轭下稳定的紧子集。考察 S 的邻域中取复值的全纯函数的芽所成的 C-代数 \(\mathcal{O}(S)\)，赋以 I, p.~49 第 1 小节中定义的复局部凸空间结构。若 U 是 S 在 C 中的开邻域，\(h: U \to C\) 是全纯函数，则 U 在复共轭下的像 V 是 S 在 C 中的开邻域，且 \(h^{*}: w \mapsto \overline{h(\overline{w})}\) 是 V 上的全纯函数。过渡到归纳极限，就得到一个连续对合 \(f\mapsto f^{*}\)（在代数 \(\mathcal{O}(\mathrm{S})\) 上）。特别地，我们有：

\[
(f + g) ^ {*} = f ^ {*} + g ^ {*} \quad (f g) ^ {*} = f ^ {*} g ^ {*} \quad (\lambda f) ^ {*} = \bar {\lambda} f ^ {*}
\]

对 \(f,g\)（属于 \(\mathcal{O}(\mathrm{S})\)）与 \(\lambda\)（属于 \(\mathbf{C}\)）成立。

我们用 \(\mathcal{O}_{\mathbf{R}}(\mathrm{S})\) 表示由芽 \(f \in \mathcal{O}(\mathrm{S})\)（满足 \(f = f^{*}\)）组成的集合。它是一个闭全子 \(\mathbf{R}\)-代数（\(\mathcal{O}(\mathrm{S})\) 中的）。

命题 19. --- 用 z 表示 \(\mathcal{O}(\mathrm{S})\) 中的、\(\mathbf{C}\) 的恒等映射的芽。则 \(\mathcal{O}_{\mathbf{R}}(\mathrm{S})\) 是最小的闭全子 \(\mathbf{R}\)-代数（在 \(\mathcal{O}(\mathrm{S})\) 中），且包含 \(z\)。

我们有 \(z^{*} = z\)，因而 z 属于 \(\mathcal{O}_{\mathbf{R}}(\mathrm{S})\)。设 B 是 \(\mathcal{O}(\mathrm{S})\) 的包含 z 的闭全子 R-代数。映射 \(f \mapsto f + f^{*}\)（从 \(\mathcal{O}(\mathrm{S})\) 到 \(\mathcal{O}_{\mathbf{R}}(\mathrm{S})\) 中）是连续且满射的，并且 S 的邻域中的全纯有理函数的芽的集合在 \(\mathcal{O}(\mathrm{S})\) 中稠密（I, p.~69 的定理 3）。因此，要证明 B 包含 \(\mathcal{O}_{\mathbf{R}}(\mathrm{S})\)，只需证明：若 f 是这样一个有理函数的芽，则有 \(f + f^{*} \in B\)。

存在 \(\mathbf{C}[X]\) 中的多项式 P 与 Q，使得 Q 在 S 的任何点上都不为零，并且有 \(f = \frac{\mathrm{P}(z)}{\mathrm{Q}(z)}\)。用 \(P^{*}\) 与 \(Q^{*}\) 表示把 P 与 Q 的系数换成其共轭而得到的多项式。于是有 \(\mathrm{P}(z)^{*} = \mathrm{P}^{*}(z)\) 与 \(\mathrm{Q}(z)^{*} = \mathrm{Q}^{*}(z)\)。由于 S 在复共轭下稳定，多项式 \(Q^{*}\) 在 S 的任何点都不为零。因此芽 \(\mathrm{Q}^{*}(z)\) 与 \((\mathrm{QQ}^{*})(z)\) 在 \(\mathcal{O}(\mathrm{S})\) 中可逆，并且

\[
f + f ^ {*} = \frac {\mathrm{P} (z)}{\mathrm{Q} (z)} + \frac {\mathrm{P} ^ {*} (z)}{\mathrm{Q} ^ {*} (z)} = \frac {(\mathrm{PQ} ^ {*} + \mathrm{P} ^ {*} \mathrm{Q}) (z)}{(\mathrm{QQ} ^ {*}) (z)}.
\]

由于多项式 PQ* + P*Q 与 QQ* 具有实系数，而 B 是 \(\mathcal{O}(\mathrm{S})\) 的包含 \(z\) 的全子 R-代数，故元素 \(f + f^{*}\) 属于 B。命题证明完毕。

设 A 是 R 上的一个完备可赋范幺代数。设 x 是 A 的元素。元素 \(1 \otimes x\)（属于代数 \(\mathrm{A}_{(\mathbf{C})}\)）的谱称为 x 的复谱，记作 \(\mathrm{Sp}_{\mathrm{A}_{(\mathbf{C})}}(x)\)。它与集合 R 的交正是 x 相对于 A 的谱 \(\mathrm{Sp}_{\mathrm{A}}(x)\)，后者有时称为 x 的实谱。复谱 \(\mathrm{Sp}_{\mathrm{A}_{(\mathbf{C})}}(x)\) 是 C 的紧子集，在复共轭下稳定。当代数 A 不化为 0 时它非空。

设 \(x\) 是 A 的元素。\(1 \otimes x \in \mathrm{A}_{(\mathbf{C})}\) 的谱半径等于谱半径 \(\varrho(x)\)（即 \(x\) 的谱半径）。它是最小的实数 \(r \geqslant 0\)，使得 \(|\lambda| \leqslant r\)（对所有 \(\lambda \in \mathrm{Sp}_{\mathrm{A}_{(\mathbf{C})}}(x)\)）。我们有

\[
\varrho (x) = \lim _ {n \to + \infty} \| x ^ {n} \| ^ {1 / n} = \inf _ {n > 0} \| x ^ {n} \| ^ {1 / n}
\]

对 A 上定义其拓扑的每个范数成立。事实上，可以假设 A 上的范数是 \(\mathrm{A}_{(\mathbf{C})}\) 上某个范数（它定义 \(\mathrm{A}_{(\mathbf{C})}\) 的拓扑）在 A 上的限制，并应用 I, p.~20 的命题 1。

我们用 \(u \mapsto \overline{u}\) 表示 R-代数 \(A_{(C)}\) 的把 \(\lambda \otimes a\) 映为 \(\overline{\lambda} \otimes a\) 的自同态。它是连续的。

引理. --- 对每个 \(f \in \mathcal{O}(\mathrm{Sp}_{\mathrm{A}_{(\mathbf{C})}}(x))\)，有 \(f^{*}(1 \otimes x) = \overline{f(1 \otimes x)}\)。

映射 \(f \mapsto f(1 \otimes x)\) 与 \(f \mapsto \overline{f^*(1 \otimes x)}\) 都是连续保单位元的 \(\mathbf{C}\)-代数同态：从 \(\mathcal{O}(\mathrm{Sp}(x))\) 到 \(\mathrm{A}_{(\mathbf{C})}\) 中，且都把 \(z\) 映为 \(1 \otimes x\)；因此它们相等（I, p.~74, 定理 5）。

命题 20. --- 对每个 \(f \in \mathcal{O}_{\mathbf{R}}(\mathrm{Sp}_{\mathrm{A}_{(\mathbf{C})}}(x))\)，存在 A 的唯一的元素 \(f(x)\)，使得 \(f(1 \otimes x) = 1 \otimes f(x)\)（在 \(\mathrm{A}_{(\mathbf{C})}\) 中）。映射 \(f \mapsto f(x)\)（从 \(\mathcal{O}_{\mathbf{R}}(\mathrm{Sp}_{\mathrm{A}_{(\mathbf{C})}}(x))\) 到 A 中）是唯一的连续保单位元 R-代数同态，它把 C 的恒等映射在 \(\mathcal{O}_{\mathbf{R}}(\mathrm{Sp}_{\mathrm{A}_{(\mathbf{C})}}(x))\) 中的芽映为 \(x\)。

我们记 \(S = \text{Sp}_{\text{A}(\mathbf{C})}(x)\)。由上面的引理，对每个芽 \(f \in \mathcal{O}_{\mathbf{R}}(\text{Sp}(x))\)，有 \(f(1 \otimes x) = \overline{f(1 \otimes x)}\)。由此得出第一个断言。用 \(z\) 表示 \(\mathcal{O}_{\mathbf{R}}(\text{S})\) 中的、\(\mathbf{C}\) 的恒等映射的芽。映射 \(f \mapsto f(x)\) 是从 \(\mathbf{R}\)-代数 \(\mathcal{O}_{\mathbf{R}}(\text{Sp}(x))\) 到 \(A\) 中的连续保单位元同态，它把 \(z\) 映为 \(x\)。由命题 19，这是唯一的这样一个同态，因为任何具有这些性质的态射在每个子 \(\mathbf{R}\)-代数（属于 \(\mathcal{O}(\text{S})\) 且包含 \(z\)）上都唯一确定。

设 \(f \in \mathcal{O}_{\mathbf{R}}(\mathrm{Sp}_{\mathrm{A}(\mathbf{C})}(x))\)。元素 \(f(x)\) 属于 A 的每个包含 \(x\) 的闭全子代数（命题 19），因而属于 \(x\) 在 A 中的双交换子。\(f(x)\) 的复谱等于 \(f(\mathrm{Sp}(x))\)（I, p.~75, 命题 8）。对所有 \(g \in \mathcal{O}_{\mathbf{R}}(f(\mathrm{Sp}_{\mathrm{A}(\mathbf{C})}(x)))\)，有 \(g \circ f \in \mathcal{O}_{\mathbf{R}}(\mathrm{Sp}_{\mathrm{A}(\mathbf{C})}(x))\) 且（同前）\(g \circ f(x) = g(f(x))\)。

设 U 是 C 的在复共轭下稳定的开子集。A 中复谱包含在 U 中的元素 x 组成的集合 \(\Omega\) 在 A 中是开的（I, p.~76, 命题 10）。设 f 是 U 上满足 \(f^{*} = f\) 的全纯函数。映射 \(x \mapsto f(x)\)（从 \(\Omega\) 到 A 中）是解析的（同前）。

设 A、B 是 R 上的完备可赋范幺结合代数，\(\varphi\colon A\to B\) 是连续的保单位元代数态射。设 \(x\in A\)。\(\varphi(x)\) 的复谱包含在 x 的复谱中，并且对每个 \(f\in\mathcal{O}_{\mathbf{R}}(\mathrm{Sp}_{\mathrm{A}(\mathbf{C})}(x))\)，有 \(f(\varphi(x))=\varphi(f(x))\)。这立即由复情形的相应命题得出（I, p.~75, 命题 8）。

